\section{Experimental Setup}\label{sec:setup}

\textit{Research questions.}
\textbf{RQ1:} How often do translations from each converter deploy, pass the lint, and agree with the
reference?
\textbf{RQ2:} What share of translations passes G1 and G2 but fails G3?
\textbf{RQ3:} Which categories remain hard, and where do the remaining failures come from?
\textbf{RQ4:} How do model size and quantization affect agreement and latency on local hardware?
\textbf{RQ5:} Which reconciliation checks detect injected data faults?

\textit{Platform.} All experiments ran on one Mac Studio (\HwChip) with \HwMemory{} of unified memory
(\HwCores{} cores) under macOS \HwMacOS. SQL Server 2022 Developer Edition (build \SqlBuild) ran in Docker
\DockerVersion{} as the \texttt{amd64} image \texttt{mcr.microsoft.com/mssql/server} (digest \texttt{\ImageDigest}) under
Rosetta emulation. Models were served by Ollama~\cite{ollama}, which
uses llama.cpp~\cite{llamacpp}. We evaluate \NumModels{} configurations from two model
families~\cite{hui2024qwen,deepseek2024v2}: Qwen2.5-Coder at 7B, 14B, and 32B with 4-bit quantization
(Q4\_K\_M), the 32B model with 8-bit quantization (Q8\_0), and DeepSeek-Coder-V2 Lite, a 16B
mixture-of-experts model with 4-bit quantization (Q4\_0). Model digests are recorded in the artifact.
\maybeinput{tables/hosted_setup.tex}

\textit{Converters.} The rule baseline applies 22 regular-expression rewrites and normalises quotes; it
covers the lint-visible categories and deliberately does not restructure joins or grouping, so it is a
lower bound rather than a substitute for a commercial converter. The plain prompt asks for a translation
that preserves behaviour. The checklist prompt adds one line per G2 rule and three lines for silent
categories that the lint cannot detect. The repair loop runs Algorithm~\ref{alg:agent} with $K=3$.
Temperature is 0.2 with fixed seeds, and each LLM configuration runs three times per case; the rule baseline
is deterministic and runs once. The held-out set therefore has 276 graded runs: 6 procedures, each
translated by 5 models with 3 LLM converters 3 times (270), plus 6 rule-baseline runs.

\textit{Prompts.} Both prompts share the instruction below (shortened); the full text is in the artifact.
\begin{lstlisting}
You are an expert database migration engineer.
Translate SAP ASE (Sybase) Transact-SQL into
Microsoft SQL Server 2022 Transact-SQL. Preserve
the runtime behaviour of the original exactly
(same result sets, same NULL handling, same
errors raised). Keep object names unchanged.
Return only the translated code in a single
sql block.
\end{lstlisting}

\textit{Metrics.} G1 and G2 rates are computed over all cases and G3 agreement over cases scored
differentially. The \emph{silent} rate is the share of differential runs that pass G1 and G2 but fail G3.
The three seeds gave the same G3 outcome in \RepAgreement\% of case and configuration pairs, so we treat the
case, not the run, as the unit of analysis: intervals are 95\% bootstrap confidence intervals (CIs) over
cases (10{,}000 resamples)~\cite{efron1993}, and paired comparisons use the exact McNemar
test~\cite{mcnemar1947} on per-case majority outcomes. Latency is the wall-clock time of the model calls
for one case; for the repair loop it also includes the G1 and G2 checks between attempts, but never G3.

\textit{Model selection.} The ``strongest model'' is chosen from the results, as the configuration with the
highest agreement for the repair loop. Statistical tests on that model are therefore exploratory. To show
that the main result does not depend on this choice, we also test the checklist effect for every model,
with Holm correction across models.

\textit{Test adequacy.} To measure how well the inputs detect behaviour changes, we applied 23 mutation
operators to the reference translations, each making one small change of the kind seen in failed
translations (for example \texttt{IS NULL} to \texttt{= NULL}, a left join to an inner join, removing an
\texttt{ISNULL}, widening a comparison, narrowing a type, changing an error message, or dropping a
\texttt{COMMIT}). This gives \MutTotal{} mutants. A mutant is detected if it fails G1 or G3 against the
unchanged reference; we inspected every mutant that was not detected.

\textit{G4 protocol.} A seeded 2{,}000-row table with Unicode names, names containing the separator
character, 4-decimal money values, microsecond timestamps, and nullable text values of up to 6{,}000
characters is loaded as the source. For each fault type and each of \RecTrials{} trials, the target
receives one injected fault. Most faults change one row; swapping amounts between two rows changes two.
Four faults are aimed at the naive encoding: a NULL replaced by the text \texttt{<NULL>}, a value moved
across the separator between two columns, a change at the end of a value longer than 4{,}000 characters,
and one row's key replaced by another row's key. Clean trials measure false positives. Source and target
are both on SQL Server, so the study measures detector sensitivity, not cross-engine extraction.
